\section*{Exercises}
\addcontentsline{toc}{section}{Exercises}

The solutions to the exercises that do not ask for code are in
\hyperref[ch:losningar]{appendix~A}. The code is published in the course repository after the
lecture.

\exercise{Routing}{Reasoning}
\label{ex:3:1}
\begin{enumerate}[a)]
  \item Does the routing happen in the bus or in the node? Justify your answer with what a
    physical wire can and cannot do.
  \item A node receives a frame with a correct checksum, but \code{DST} is not its address. What
    should it do?
  \item Why should it \emph{not} send a \code{Nack} in that situation?
\end{enumerate}

\exercise{The error model}{Reasoning}
\label{ex:3:2}
For each of the errors --- dropped byte, flipped bit, delay --- answer:
\begin{enumerate}[a)]
  \item Does the frame end up structurally correct?
  \item Which mechanism detects the error?
  \item Which state is the parser in afterwards?
\end{enumerate}
One of the three errors should cause no problem at all. Which one, and what does it mean if it
does anyway?

\exercise{A lost acknowledgement}{Decoding}
\label{ex:3:3}
Node A sends a \code{StatusRequest} with \code{SEQ 0x0007} to node B. B receives it, processes it,
and sends an \code{Ack}. The acknowledgement is lost.
\begin{enumerate}[a)]
  \item What does A do next, and after what?
  \item What does B see then?
  \item What should B do, and what should B not do?
  \item What would have happened if B had stayed silent instead?
\end{enumerate}

\exercise{The timeout value}{Reasoning}
\label{ex:3:4}
\begin{enumerate}[a)]
  \item What happens if the timeout is too short? Give two consequences.
  \item What happens if it is too long?
  \item What information would be needed to choose the value correctly, and why does a node
    rarely have it?
\end{enumerate}

\exercise{The chain}{Reasoning}
\label{ex:3:5}
List the five mechanisms in the reliability chain, in order. For each one: which problem does it
solve, and which new problem does it create?
